\section{Baseline learning curves: dropout 0.1}
\label{sec:baseline}
Appendix Figure~\ref{fig:baseline} shows 13 recorded evaluations for each model: before training, roughly every ten minutes, and at the end. Training loss is measured on fixed training windows with dropout turned off. It is not the loss from a single training microbatch. The lower plots leave out the initial evaluation and use the same scale to make later changes easier to compare. This baseline comparison was used to decide what to investigate next.

\begin{table}[H]
\centering
\caption{Baseline loss and best recorded time. All values are sampled cross-entropy in nats/token.}
\begin{tabular}{lrrrr}
\toprule Model & Best val. & Best time (min) & Final train & Final val.\\\midrule
Small & \textbf{1.7652} & 110.1 & 1.5603 & 1.7659\\
Medium & 1.8018 & 40.1 & 1.3299 & 1.8892\\
Large & 1.8122 & 50.3 & 1.2865 & 1.9224\\\bottomrule
\end{tabular}
\end{table}

All models improve quickly during the first ten minutes. Small then improves more slowly and levels off near the end. Medium and large reach their best validation losses earlier. After that, their training losses keep falling while their validation losses rise. Large has the lowest final training loss but the highest final validation loss. This shows that fitting the training pieces well does not necessarily mean predicting unseen pieces well.

Using the best recorded validation checkpoint, small ranks first, followed by medium and large. The differences are fairly small: medium and large have best-checkpoint perplexities about 3.7\% and 4.8\% higher than small. This does not mean that more parameters always hurt. It means that larger models did not perform better with this dataset, implementation, schedule, and time limit.
